\chapter{Lifecycle: from main to exit}
\label{ch:lifecycle}

\begin{aside}
A \maya{} app starts with \code{main} and ends with the
terminal restored. Between those two moments, a precise sequence
of events happens. This chapter traces them: who runs, in what
order, why each phase matters. The mental model lets you
predict where a hang, a crash, or a hot-reload should be hooked.
\end{aside}

\section{Phase 0: pre-main}

Before \code{main} runs, C++ does work:

\begin{enumerate}[leftmargin=*]
  \item Static-storage-duration objects are initialised in
    definition order within a translation unit; order between
    TUs is unspecified.
  \item \code{constinit} variables are guaranteed to be ready
    by the start of \code{main}.
  \item Function-local static objects are deferred to first use.
\end{enumerate}

\maya{} uses \code{constinit} for theme defaults and palette
tables to avoid the static-init order fiasco. Anything that
needs to be ready before \code{main} should be \code{constinit}.

\section{Phase 1: argument parsing and config}

\code{main} typically:

\begin{lstlisting}
int main(int argc, char** argv) {
    auto cfg = build_config(argc, argv);
    if (cfg.show_help) { print_help(); return 0; }
    return maya::run<App>(cfg);
}
\end{lstlisting}

Three steps:

\begin{itemize}[leftmargin=*]
  \item Parse args.
  \item Decide whether to run (or print help / version).
  \item Call \code{run<App>}.
\end{itemize}

The handoff to \code{run} is where \maya{}'s lifecycle starts.

\section{Phase 2: run setup}

\code{run<P>(cfg)} does the following before the loop starts:

\begin{enumerate}[leftmargin=*]
  \item Construct the \code{Context} (Chapter~\ref{ch:runtime}).
  \item Open the terminal and enter raw mode.
  \item Install signal handlers (SIGINT, SIGWINCH, SIGTERM).
  \item Start the worker pool.
  \item Construct the renderer, the event loop, the interpreter.
  \item Call \code{P::init()} to get the initial model and
    command.
  \item Interpret the initial command (e.g. fire a startup
    task).
\end{enumerate}

If any step fails, the runtime unwinds the prior steps and
returns an error code without entering the main loop. The
terminal is restored, the workers are joined, and resources are
released.

\subsection*{Failure modes during setup}

Common failures:

\begin{itemize}[leftmargin=*]
  \item Terminal does not exist (running under cron, no tty).
  \item Permission denied (the user's stdin is read-only).
  \item Raw mode setup fails (rare; usually permission).
  \item Worker pool cannot start (out of threads).
\end{itemize}

\maya{} propagates these as exit codes and writes a single
line to stderr explaining.

\section{Phase 3: the main loop}

The loop from Chapter~\ref{ch:runtime}:

\begin{enumerate}[leftmargin=*]
  \item Render the current model.
  \item Compute the current subscription set.
  \item Wait for an event.
  \item Translate the event into a message.
  \item Run \code{update}; replace the model.
  \item Interpret the resulting command.
  \item Repeat.
\end{enumerate}

The loop ends when an interpretation sets the running flag to
false. The only message that does this is \code{Cmd::quit}; the
typical trigger is the user pressing \code{q} or \code{Ctrl-C}.

\subsection*{What can interrupt the loop}

\begin{itemize}[leftmargin=*]
  \item Quit command (graceful exit).
  \item SIGINT / SIGTERM (signal handler sets running flag).
  \item An exception escaping the runtime (rare; logged and
    aborts).
  \item Out-of-memory (terminates).
\end{itemize}

\section{Phase 4: shutdown}

When the loop exits, destructors run in reverse construction
order:

\begin{enumerate}[leftmargin=*]
  \item \code{Interpreter}: drain pending tasks, signal workers
    to stop, join.
  \item \code{Renderer}: flush any pending bytes, reset cursor,
    restore default style.
  \item \code{EventLoop}: close epoll/kqueue/handle, clean up
    the wake fd.
  \item \code{Terminal}: leave raw mode, leave alt screen if
    entered.
  \item \code{Context}: close logger, release config.
\end{enumerate}

After destructors, \code{run} returns. The exit code is the
last value the application set (typically 0).

\section{The order matters}

If you reverse the order ---  for example, restore the terminal
before joining workers --- a worker thread might still be
writing bytes when the terminal is back in cooked mode, and the
user sees corrupted output.

The construction order in \code{run} sets the destruction
order. The pattern: construct in dependency order; let RAII
handle the reverse.

\subsection*{What about atexit?}

\code{atexit} handlers run after \code{main} returns. \maya{}
registers a fallback handler that calls \code{leave\_raw\_mode}
in case the normal destructor did not. This handles
\code{std::exit} and uncaught exceptions.

\section{Hot reload}

In development, a hot-reload pattern is sometimes useful: detect
a code change, rebuild, restart the app while preserving the
model.

The pattern:

\begin{enumerate}[leftmargin=*]
  \item The app writes its model to a file on \code{Quit}.
  \item A watcher process notices the executable changed.
  \item The watcher kills the app, rebuilds, restarts it.
  \item The new app reads the model file and restores state.
\end{enumerate}

\maya{}'s architecture makes this tractable: the model is a
plain struct, serialisable to JSON or another format.

\section{Daemonisation}

\maya{} apps are interactive; daemonising them does not make
sense. If you want background work, use \code{Cmd::task}; if
you want a long-running service, write a regular service, not
a TUI.

\section{Subprocess management}

A TUI app sometimes runs subprocesses (git, ripgrep, an LLM).
The cleanest pattern:

\begin{lstlisting}
[&](RunGitDiff) {
    return std::pair{m, Cmd<Msg>::task([] -> Msg {
        auto out = exec(\"git\", {\"diff\", \"--stat\"});
        return GitDiffResult{std::move(out)};
    })};
}
\end{lstlisting}

The subprocess runs on a worker thread; the result message
arrives back. The user does not see a frozen UI.

\subsection*{Pty subprocesses}

For interactive subprocesses (a shell, vim), you need a
pseudo-terminal. \maya{} does not ship a pty manager, but the
patterns are well-known (forkpty on POSIX, ConPTY on Win32).
Wrap the pty in a \code{Cmd::task} that pumps bytes back as
messages.

\section{Crash recovery}

If the app crashes during a long operation, the user wants:

\begin{itemize}[leftmargin=*]
  \item The terminal restored (so they can keep using it).
  \item A crash log written to disk.
  \item State restorable on restart.
\end{itemize}

The signal handler from Chapter~\ref{ch:observability} handles
the first two. The third requires periodic state saves
(Chapter~\ref{ch:config}).

\section{Exit codes}

Convention:

\begin{itemize}[leftmargin=*]
  \item \code{0}: normal exit.
  \item \code{1}: user-facing error.
  \item \code{2}: internal error (bug).
  \item \code{3}: configuration error.
  \item \code{130}: interrupted by SIGINT (128 + 2).
\end{itemize}

\maya{} returns the value from \code{run<P>(cfg)}; you control
the rest via \code{Cmd::quit\_with(code)}.

\section{Pitfalls}

\begin{warning}
\textbf{Static destructors after main.} If you have a static
object whose destructor writes to stdout, the writes happen
\emph{after} \code{main} returns. By then \maya{} has restored
the terminal but the program is still alive. The writes corrupt
the user's prompt. Avoid output-producing destructors in
statics.
\end{warning}

\begin{warning}
\textbf{Forgetting to set a quit hotkey.} An app with no
\code{q} or \code{Ctrl-C} handler can only be exited by the
user killing the process. Always provide a clean exit path.
\end{warning}

\begin{warning}
\textbf{Holding the terminal during a long task.} If
\code{update} takes seconds, the UI freezes and the user
suspects a hang. Move long tasks to \code{Cmd::task}; show a
spinner while they run.
\end{warning}

\begin{warning}
\textbf{Joining workers in main thread before they finish.}
The interpreter's destructor waits for workers to drain.
If a worker is stuck, the shutdown hangs. Use timeouts on
tasks that might block forever.
\end{warning}

\section*{Workshop}

\begin{enumerate}[leftmargin=*]
  \item Add startup logging to a small app: log each phase as
    it starts. Verify the order matches this chapter.
  \item Force a crash in \code{update}. Confirm the terminal
    is still usable after the process dies.
  \item Add a hot-reload cycle: persist model on quit; restore
    on startup; observe the seamless restart.
\end{enumerate}

\begin{recap}
A \maya{} app's lifecycle has four phases: pre-main static
init, run setup, main loop, shutdown. Each phase has clear
ordering and clear failure modes. RAII handles shutdown order
correctly. The architecture invites long-running tasks to be
\code{Cmd::task}s and signals to be subscriptions; the loop
itself is small and uniform.
\end{recap}
